\chapter{ОБЗОР ЛИТЕРАТУРЫ}

Настоящая глава охватывает три направления, на пересечении которых находится данное исследование: сначала рассматриваются геномные языковые модели, затем методы токенизации ДНК и после этого - бенчмарки для оценки их качества.

\section{Языковые модели ДНК}

Языковые модели ДНК обучаются на нуклеотидных последовательностях в режиме самосупервизии и формируют представления, которые затем можно использовать для задач геномики~\cite{benegasGenomicLanguageModels2024}. Геном человека содержит порядка $3{,}1 \times 10^9$ пар оснований~\cite{nurgalievT2TCHm13CompleteHuman2022}, причём значительная часть функциональной информации скрыта в некодирующих регионах, которые сложно аннотировать напрямую. Языковые модели позволяют извлекать эту информацию из данных без явных меток.

\subsection{Архитектуры на основе Transformer}

Transformer~\cite{vaswaniAttentionAllYou2017}, изначально созданный для NLP, довольно быстро перекочевал в геномику.

DNABERT~\cite{jiDNABERTPretrainedBidirectional2021} - одна из первых моделей, применивших BERT к ДНК. Она работает с перекрывающимися $k$-мерами ($k = 3, \ldots, 6$) в режиме MLM. Проблема в том, что фиксированные $k$-меры привязывают модель к одному масштабу, а словарь растёт как $4^k$. DNABERT-2~\cite{zhouDNABERT2EfficientFoundation2024} заменяет $k$-меры на BPE, что даёт токены переменной длины и позволяет обучаться на мультивидовых данных.

Nucleotide Transformer~\cite{dalla-torreNucleotideTransformerBuilding2025} масштабирует подход до $2{,}5 \times 10^9$ параметров, обучаясь на последовательностях более чем 3000 видов с неперекрывающимися 6-мерами. При этом масштаб помогает на downstream-задачах, но фиксированные 6-меры остаются ограничением: границы $k$-меров не совпадают с функциональными границами в геноме.

GROVER~\cite{groverLearnsSequenceContext2024} обучается на человеческом геноме с BPE-токенизацией и посимвольной инициализацией. Важно, что выученный BPE-словарь частично отражает биологически значимые подстроки, хотя сам BPE оптимизирует частотность, а не биологию.

\subsection{Модели на основе State Space Models}

State Space Models (SSM)~\cite{guMambaLinearTimeSequence2023} - альтернатива механизму внимания для длинных последовательностей. Здесь ключевое преимущество - линейная сложность $O(L)$ вместо $O(L^2)$ у стандартного Transformer.

HyenaDNA~\cite{nguenHyenaDNALongRangeGenomic2024} построен на архитектуре Hyena (длинные свёртки) и работает посимвольно с окном до $10^6$ нуклеотидов. Посимвольная обработка снимает проблему произвольных границ, но ценой очень длинных последовательностей, для обработки которых нужна специализированная архитектура.

Caduceus~\cite{schiffCaduceusBiDirectionalEquivariant2024} добавляет к SSM два свойства: двунаправленность и RC-эквивариантность. ДНК - двуцепочечная молекула, и многие функциональные элементы (мотивы TF, промоторы) работают на обеих цепях. RC-эквивариантная архитектура гарантирует, что представления прямой и комплементарной цепей связаны фиксированным преобразованием. Это сильный индуктивный bias для геномных данных.

Evo~\cite{nguyenSequenceModelingDesign2024} и Evo~2~\cite{nguyenEvo2GenomeModeling2026} - масштабирование SSM до $7 \times 10^9$ параметров на геномах всех доменов жизни. Evo~2 использует StripedHyena с окном до $1{,}6 \times 10^5$ нуклеотидов и генерирует синтетические последовательности с биологически правдоподобными свойствами.

\subsection{Гибридные модели}

HybriDNA~\cite{maHybriDNAHybridTransformerMamba22025} совмещает блоки Transformer и Mamba-2~\cite{daoTransformersAreSSMs2024}, получая выразительность attention и эффективность SSM в одной архитектуре. Модель посимвольная, и при меньшем числе параметров достигает конкурентных результатов.

AlphaGenome~\cite{avsecAlphaGenomeAdvancingRegulatory2025} от DeepMind ориентирован на предсказание эффектов регуляторных вариантов - от уровня нуклеотидов до хромосомного масштаба.

\subsection{Общие тенденции}

Если собрать обзор воедино, можно выделить несколько наблюдений.

Контекстное окно постоянно растёт: от $512$ нуклеотидов в ранних DNABERT до $10^6$ в HyenaDNA. Это логично: взаимодействия вроде энхансер--промотор~\cite{schoenfelderLongrangeEnhancerPromoter2019} работают на расстояниях в десятки тысяч пар оснований.

Во многих сильных результатах (например, HyenaDNA, Caduceus, Evo) используется посимвольная токенизация, а не $k$-меры или BPE. Это указывает на ограничения статических методов в части универсальности. При этом посимвольная токенизация создаёт вычислительную проблему: длина последовательности напрямую определяет стоимость - $O(L^2)$ для Transformer или $O(L)$ для SSM, но с большим постоянным множителем.

Наконец, несколько недавних работ~\cite{lindseyImpactTokenizerSelection2025, vishniakovGenomicFoundationlessModels2024} прямо показывают, что выбор токенизации существенно влияет на качество. Иными словами, нужен подход, который сжимает последовательность, но не теряет функционально важную информацию.


\section{Методы токенизации ДНК-последовательностей}

\subsection{Посимвольная токенизация}

Самый простой вариант: каждый нуклеотид (A, T, C, G) - отдельный токен. Словарь минимален ($|V| = 4$ плюс специальные токены), артефактов разбиения нет, разрешение максимальное. Именно так работают HyenaDNA~\cite{nguenHyenaDNALongRangeGenomic2024}, Caduceus~\cite{schiffCaduceusBiDirectionalEquivariant2024} и Evo~\cite{nguyenSequenceModelingDesign2024}.

Недостаток очевиден: длина входа равна длине ДНК. Для моделей с $O(L^2)$ это проблема. Для SSM - менее критично, но затраты всё равно линейны по длине.

\subsection{Токенизация $k$-мерами}

Разбиение на подстроки длины $k$ сокращает вход в $k$ раз (при неперекрывающихся $k$-мерах). Использовалось в DNABERT ($k = 3, \ldots, 6$)~\cite{jiDNABERTPretrainedBidirectional2021} и Nucleotide Transformer ($k = 6$)~\cite{dalla-torreNucleotideTransformerBuilding2025}.

Словарь растёт экспоненциально: $|V| = 4^k$. При $k = 6$ это 4096 токенов - управляемо. При $k \geq 8$ уже $> 65\,000$, и embedding matrix становится неоправданно большой. Но главная проблема не в этом: границы $k$-меров произвольны. Один и тот же мотив может быть разбит по-разному в зависимости от сдвига начала разбиения.

\subsection{BPE-токенизация}

BPE~\cite{sennrichNeuralMachineTranslation2016} строит словарь итеративным слиянием наиболее частотных пар символов. Используется в DNABERT-2~\cite{zhouDNABERT2EfficientFoundation2024} и GROVER~\cite{groverLearnsSequenceContext2024}. Токены получаются переменной длины, что частично решает проблему фиксированного масштаба.

Однако BPE оптимизирует статистическое сжатие, а не биологическую информативность. Повторяющиеся элементы генома (Alu, LINE, SINE - около 45\% генома человека~\cite{landerInitialSequencingAnalysis2001}) получают компактные токены благодаря высокой частотности. Уникальные регуляторные элементы остаются фрагментированными. Кроме того, BPE-словарь строится до обучения и остаётся фиксированным.

\subsection{Ограничения статических подходов}

Общая проблема всех трёх подходов - правила разбиения не зависят от контекста. Между тем геном неоднороден: плотность функциональной информации в промоторе и в межгенном повторе отличается на порядки, а статическая токенизация применяет к ним одни и те же правила.

Lindsey~и~др.~\cite{lindseyImpactTokenizerSelection2025} систематически сравнивают токенизаторы и показывают, что ни один статический метод не доминирует на всех задачах. Eapen~\cite{eapenGenomicTokenizerBiologydriven2025} предлагает биологически обоснованные правила сегментации, но они определяются вручную и не обучаются вместе с моделью.


\section{Адаптивная и обучаемая токенизация}

\subsection{Адаптивная токенизация в NLP}

Проблема фиксированной токенизации активно исследуется и в NLP. Для данной работы особенно релевантны четыре подхода.

Charformer~\cite{tayCharformerFastCharacter2022} вводит Gradient-Based Subword Tokenization (GBST) - модуль мягкой токенизации, обучаемый совместно с моделью. GBST вычисляет взвешенную комбинацию группировок символов через обучаемые скоры. Ограничение - одноуровневая группировка без ограничения непересечения, из-за чего границы токенов трудно интерпретировать.

Funnel Transformer~\cite{daiFunnelTransformerFiltering2020} применяет stride pooling для постепенного сокращения длины. Подход эффективен, но не адаптивен: коэффициент сжатия одинаков для всех частей входа.

MegaByte~\cite{yuMegaBytePredictingMillionbyte2023} разделяет обработку на локальный и глобальный Transformer с фиксированным patching на байтовом уровне. Это позволяет обрабатывать до $10^6$ байт, но границы patch'ей не адаптируются к содержимому.

Byte Latent Transformer (BLT)~\cite{bltByteLatentTransformer2024} от Meta FAIR ближе всего к идее адаптивности. BLT определяет границы patch'ей по энтропии предсказания следующего байта: области с высокой неопределённостью получают более мелкое разбиение. В экспериментах авторов на сопоставимом вычислительном бюджете адаптивные границы оказываются эффективнее фиксированных.

Однако ни один из этих подходов не использует предметно-специфические сигналы: все они разработаны для текста общего назначения.

\subsection{Адаптивная токенизация ДНК}

В геномике ситуация развивалась немного иначе: в 2024--2026~годах появилось несколько работ, предлагающих обучаемую токенизацию именно для ДНК.

MxDNA~\cite{qiaoMxDNAAdaptiveDNA2024} (NeurIPS 2024) реализует подход через разреженную смесь свёрточных экспертов с деформируемыми свёртками. Выбор эксперта определяет масштаб анализа. Авторы сообщают об улучшениях на ряде задач по сравнению с фиксированными токенизаторами, но механизм выбора не использует биологические приоры и не даёт явного контроля над границами.

MergeDNA~\cite{mergeDNAContextAware2026} вычисляет попарные оценки сходства между соседними токенами и сливает наиболее похожие пары. По механизму слияния это ближе всего к предлагаемому в данной работе подходу, но MergeDNA - одноуровневый, без рекурсии и без дифференцируемого DP.

DNAChunker~\cite{dnachunkerLearnableTokenization2026} обучает модуль сегментации на чанки переменной длины совместно с языковой моделью. Концептуально близок к обучаемой сегментации, но без рекурсии и без биоприоров.

EvoLen~\cite{evoLenEvolutionGuided2026} - одна из первых работ, явно использующих биологические приоры для управления гранулярностью: длина токена зависит от оценки эволюционной консервативности. Высоко консервативные регионы получают более короткие токены. Правда, правила применяются детерминистически, без обучаемого компонента.

DNAMotifTokenizer~\cite{dnaMotifTokenizer2025} сегментирует по известным мотивам TF из базы JASPAR~\cite{rauluseviciuteJASPAR2024}. Токены получаются биологически интерпретируемые, но метод зависит от полноты аннотаций и не работает для организмов с плохо изученными регуляторными элементами.

BioToken/BioFM~\cite{bioTokenBioFM2025} и PatchDNA~\cite{patchDNA2026} сообщают, что даже простой учёт биологической структуры при построении токенов улучшает downstream-качество. Guided Tokenization~\cite{guidedTokenization2026} демонстрирует схожую тенденцию для задач аннотации и классификации.

\subsection{Сравнительный анализ подходов}

Чтобы не смешивать разные идеи в один список, существующие подходы удобно разложить по нескольким осям: обучаемый или детерминистический механизм определения границ, наличие биологических приоров, число уровней иерархии и способ обеспечения дифференцируемости.

В рассмотренных работах не встречается комбинация всех четырёх свойств сразу: рекурсивное многоуровневое слияние, дифференцируемое ограничение непересечения через DP, интеграция биоприоров как обучаемого индуктивного смещения и температурный annealing для минимизации разрыва обучение--инференс. Именно эту комбинацию мы предлагаем в главе~2.


\section{Бенчмарки для оценки геномных языковых моделей}

Для систематической оценки геномных моделей нужны стандартизированные бенчмарки.

Genomic Benchmarks~\cite{gresovaGenomicBenchmarksCollection2023} - коллекция задач классификации (промоторы, энхансеры, транскрипционные факторы). Задачи относительно просты и не покрывают zero-shot и probing сценарии, но удобны для базового сравнения.

BEND~\cite{tangBENDBenchmarkDNA2024} добавляет более биологически значимые задачи: структура хроматина, эффекты мутаций, функциональная аннотация. Есть задачи разной сложности и несколько режимов оценки.

Wei~и~др.~\cite{weiGenomicBenchmarksDNA2025} провели масштабное сравнение DNA foundation models на расширенном наборе задач. Результат примечателен: ни одна модель не доминирует - разные архитектуры и токенизации лучше работают на разных задачах.

Vishniakov~и~др.~\cite{vishniakovGenomicFoundationlessModels2024} пошли дальше и критически оценили саму парадигму предобучения, показав, что оно не всегда улучшает downstream-результаты. Этот результат заставляет задуматься об альтернативных архитектурных решениях - в том числе адаптивной токенизации.

DART-Eval~\cite{patel2025dartevalcomprehensivednalanguage} - на данный момент один из наиболее полных бенчмарков для DNA~LM на задачах регуляторной ДНК. Он включает пять групп задач:
\begin{enumerate}
  \item обнаружение мотивов (zero-shot) - распознавание 1443 мотивов TF через likelihood scoring;
  \item предсказание связывания TF - probing и fine-tuning;
  \item предсказание доступности хроматина;
  \item предсказание клеточно-специфичной регуляторной активности;
  \item предсказание эффектов генетических вариантов (caQTLs, dsQTLs).
\end{enumerate}

Для оценки адаптивной токенизации DART-Eval подходит по двум причинам. Сами задачи сфокусированы на регуляторной ДНК, где адаптивная токенизация должна давать наибольший выигрыш за счёт мелкозернистого разрешения на функциональных мотивах. Параллельно авторы DART-Eval показали, что существующие DNA~LM не обеспечивают устойчивого преимущества над ab~initio моделями (ChromBPNet), а значит, пространство для архитектурных улучшений действительно остаётся.


\section{Выводы по главе~1}

Итак, существующие геномные языковые модели в основном используют статическую токенизацию - посимвольную, $k$-мерную или BPE. Ни один из этих вариантов в явном виде не учитывает функциональную неоднородность генома.

Появляющиеся работы по обучаемой токенизации ДНК (MxDNA, MergeDNA, EvoLen и др.) решают отдельные аспекты проблемы, но пока не объединяют рекурсивное слияние, дифференцируемое DP и биологические приоры в единую архитектуру. Опыт NLP (BLT, Charformer) при этом показывает, что адаптивные границы нередко работают лучше фиксированных, но перенос этих идей в геномику с учётом предметных сигналов до сих пор остаётся открытым.

DART-Eval обеспечивает адекватную базу для сравнительной оценки на задачах регуляторной геномики.
